\begin{textsample}{NEG Dim 7 – human – Score -2.00 – rj/rj\_ef\_intro\_1\_p7.txt}  \label{ex:f7_neg_014}
Destaca -se, portanto, na etapa de Ensino Fundamental Anos Iniciais, os estudos de currículo que viabilizam o que acontece nas escolas, o sentido que a mesma tem, na valorização do cotidiano tendo como interlocutor Michel de Certeau.
Na visão deste autor, os conhecimentos ocorrem em diferentes contextos, para tal o respeito à produção de cada escola, o respeito à elaboração de cada Proposta Pedagógica \textbf{define} suas “ feituras ” e suas “ artes de fazer ”.
Ainda, as práticas cotidianas são tecidas por meio de táticas, contudo, compreendemos que currículo se “ produz ” e é praticado por pessoas nos espaços tempos, entende -se então que currículo se faz no cotidiano escolar, onde o que é considerado é a realidade.
De pleno acordo e entendimento sobre autonomia e poder autoral do professor que está na escola, produzindo currículos no dia a dia, não negamos as riquezas do cotidiano, tanto quanto não podemos negar a existência de mecanismos formais como as normativas que nos regem, mas podemos com isso dialogar com Certeau ( 1998, p. 99-100 ) quando trata de tática e estratégia ao citar :

% matched lemmas: definir
\end{textsample}
